\section*{第 \ref{ch:b2:population-genetics} 章——种群遗传学}

\begin{solution}{exo:b2:population-genetics:1}
$p(M) = (1280 + 320)/2000 = 0.8$，$q(N) = 0.2$。期望的 $MM$：
$0.64\times 1000 = 640$；$MN$：$2\times 0.8\times 0.2\times 1000 =
320$；$NN$：40——恰与观测数一致。该种群
在该位点处于平衡状态。
\end{solution}

\begin{solution}{exo:b2:population-genetics:2}
$q = \sqrt{1/20000} = 0.0071$；携带者 $2pq \approx 0.014$，
约每 70 人一人。携带者数量是患者的 $2p/q \approx
280$ 倍。
\end{solution}

\begin{solution}{exo:b2:population-genetics:3}
\omterm{def:b2:population-genetics:fitness}{适合度}：某\omterm{def:b2:meiosis-heredity:vocabulary}{基因型}的期望后代数相对于最优者的比值，
\omterm{def:b2:population-genetics:fitness}{选择系数} $s$：\omterm{def:b2:population-genetics:fitness}{适合度}的亏损，$w = 1 -
s$。\omterm{thm:b2:population-genetics:heritability}{遗传力} $h^{2}$：\omterm{def:b2:meiosis-heredity:vocabulary}{表型}方差中属于加性遗传的
比例，等价于子代对亲代回归的斜率。
演化动力：\omterm{def:b2:genome-diversification:mutation}{突变}、选择、漂变、迁移，以及非随机交配
（最后一项改变的是\omterm{def:b2:meiosis-heredity:vocabulary}{基因型}频率，而非\omterm{def:b2:population-genetics:frequencies}{等位基因频率}）。
\end{solution}

\begin{solution}{exo:b2:population-genetics:4}
\omterm{prop:b2:population-genetics:kinds}{定向选择}：椒花蛾的黑化现象，选择因子是在被烟灰熏黑的树干上
靠视觉捕食的鸟类。\omterm{prop:b2:population-genetics:kinds}{稳定选择}：人类出生体重，
选择因子是极小和极大婴儿的死亡率。\omterm{prop:b2:population-genetics:kinds}{平衡选择}：镰状细胞，
选择因子为疟疾（对 $AA$ 不利）和贫血（对
$SS$ 不利）共同作用。
\end{solution}

\begin{solution}{exo:b2:population-genetics:5}
$q_t = q_0/(1 + tq_0)$：频率减半需要 $1/q_0 = 50$ 代；
降至 $0.001$ 需要 $1/0.001 - 1/0.02 = 950$ 代。几乎
所有该\omterm{def:b2:meiosis-heredity:vocabulary}{等位基因}的拷贝都在\omterm{def:b2:meiosis-heredity:vocabulary}{杂合}子中，对选择不可见；
阻止患者繁殖（自然界本已如此，因为患者已死亡）
对人类一生内的频率
不产生可测量的影响。
\end{solution}

\begin{solution}{exo:b2:population-genetics:6}
$p = 0.3$：$\bar w = 0.09 + 0.42 + 0.8\times 0.49 = 0.902$；$p' =
(0.09 + 0.21)/0.902 = 0.333$；$\Delta p = +0.033$。$p = 0.9$：$\bar w
= 0.998$，$p' = 0.9018$，$\Delta p = +0.0018$。在 $p = 0.9$ 时，仅
$q^{2} = 0.01$ 的个体为 $aa$：选择所针对的\omterm{def:b2:meiosis-heredity:vocabulary}{等位基因}
被隐藏在\omterm{def:b2:meiosis-heredity:vocabulary}{杂合}子中，而 $\Delta p = spq^{2}/\bar w$
带有因子 $q^{2}$。
\end{solution}

\begin{solution}{exo:b2:population-genetics:7}
设 $s = 0.12$（针对 $AA$）、$t = 0.8$（针对 $SS$），则 $\hat q =
s/(s + t) = 0.12/0.92 = 0.13$。患病出生者 $\hat q^{2} = 0.017$，
约 60 个孩子中有一个——这正是\omterm{def:b2:meiosis-heredity:vocabulary}{杂合}子所获保护的
代价。
\end{solution}

\begin{solution}{exo:b2:population-genetics:8}
\omterm{def:b2:meiosis-heredity:vocabulary}{隐性}致死：$\hat q = \sqrt{\mu} = \sqrt{2\times 10^{-5}} =
0.0045$；患病出生者 $\hat q^{2} = \mu = 2\times 10^{-5}$，即
\num{50000} 人中一人。\omterm{def:b2:meiosis-heredity:vocabulary}{显性}致死：每个拷贝在出现的那一代即被
清除，因此患病出生者即为新发生的
\omterm{def:b2:genome-diversification:mutation}{突变}，$2\mu = 4\times 10^{-5}$（每次出生涉及两个配子），即
\num{25000} 人中一人——\omterm{def:b2:meiosis-heredity:vocabulary}{隐性}致死的\omterm{def:b2:meiosis-heredity:vocabulary}{等位基因}隐藏比例为二百比一，
而\omterm{def:b2:meiosis-heredity:vocabulary}{显性}致死则不隐藏任何东西。
\end{solution}

\begin{solution}{exo:b2:population-genetics:9}
$H_t = 0.5\,(1 - 1/100)^{t}$：10 代后为 $0.45$；50 代后
为 $0.30$；200 代后为 $0.067$。对于 $(1 - 1/2N)^{100} = 0.9$：
$100/(2N) \approx -\ln 0.9 = 0.105$，因此 $N \approx 475$，约 500
个繁殖个体。
\end{solution}

\begin{solution}{exo:b2:population-genetics:10}
20 个基因拷贝：$P(\text{缺失}) = 0.9^{20} = 0.12$。单个
拷贝在 20 个中的频率为 $0.05$，中性\omterm{def:b2:meiosis-heredity:vocabulary}{等位基因}以等于其频率的概率
固定，即 $0.05$。假设大陆上有，比如说，
\num{100000} 只鸟，单个拷贝的固定概率为 $5\times 10^{-6}$
。岛屿是稀有\omterm{def:b2:meiosis-heredity:vocabulary}{等位基因}易于丢失的地方，
而幸存下来的\omterm{def:b2:meiosis-heredity:vocabulary}{等位基因}可能取得优势：漂变既使种群
贫化，也使种群\omterm{def:b2:cell-differentiation:differentiation}{分化}。
\end{solution}

\begin{solution}{exo:b2:population-genetics:11}
$R = h^{2}S = 0.3\times 1.4\times \qty{1000}{L} = \qty{420}{L}$，每
代。若选用前 \qty{1}{\%} 的公牛，平均选择差为
$(1.4 + 2.7)/2 = 2.05$ 个标准差，于是 $R =
0.3\times 2050 = \qty{615}{L}$。选择响应减慢，原因在于选择
耗尽了加性方差——有利\omterm{def:b2:meiosis-heredity:vocabulary}{等位基因}趋于固定，
$h^{2}$ 下降——还因为对其他性状的相关改变
（繁殖力、健康）带来了代价，而仅针对产量的选择
并不能察觉。
\end{solution}

\begin{solution}{exo:b2:population-genetics:12}
选择作用于\omterm{def:b2:meiosis-heredity:vocabulary}{表型}。\omterm{def:b2:meiosis-heredity:vocabulary}{隐性}\omterm{def:b2:meiosis-heredity:vocabulary}{等位基因}处于\omterm{def:b2:meiosis-heredity:vocabulary}{杂合}子中
时不表现出\omterm{def:b2:meiosis-heredity:vocabulary}{表型}，因此选择无法将其清除，其
频率以 $1/t$ 的速度下降，越来越慢。中性\omterm{def:b2:genome-diversification:mutation}{突变}
不产生\omterm{def:b2:population-genetics:fitness}{适合度}差异，其命运完全取决于漂变，它
以 $1/2N$ 的概率固定。\omterm{thm:b2:population-genetics:heritability}{遗传力}为零的性状
可能受到强烈选择——幸存者可能与
种群差异很大——然而下一代并未改变，因为
这些差异并未遗传。选择只能作用于
对\omterm{def:b2:population-genetics:fitness}{适合度}有影响的可遗传差异；其余
一切要么依靠偶然，要么根本不演化。
\end{solution}

\begin{solution}{pb:b2:population-genetics:1}
\textbf{1.} 黑化个体比例 $1 - q^{2} = 1 - 0.999^{2} = 0.002$，
即每 500 只蛾中一只；\omterm{def:b2:meiosis-heredity:vocabulary}{杂合}子中 $M$ \omterm{def:b2:meiosis-heredity:vocabulary}{等位基因}的比例为
$2pq/(2pq + 2p^{2}) = q = 0.999$，几乎全部。
\textbf{2.} 取 $\bar w = 1 - sq^{2}$，淡色\omterm{def:b2:meiosis-heredity:vocabulary}{等位基因}的拷贝
在选择后为 $pq\cdot 1 + q^{2}(1 - s)$ 除以 $\bar w$，因此 $q'
= (pq + q^{2} - sq^{2})/(1 - sq^{2}) = q(1 - sq)/(1 - sq^{2})$。
\textbf{3.} $q_1 = 0.999\times(1 - 0.2997)/(1 - 0.2994) = 0.99857$：
变化量为 $-0.0014$。
\textbf{4.} $q = \sqrt{0.1} = 0.32$。
\textbf{5.} 迭代递推式，取 $s =
0.3$ 时需 28 代；近似式给出同一量级（变化量
开始时为 $-0.0003$，随 $p$ 增大而加速）。观测到的 50
代对应较小的 $s$，约为 $0.2$（递推式
对 $s = 0.2$ 给出 45 代）：选择强度为 \qtyrange{20}{30}{\%}
每代，按多数\omterm{def:b2:meiosis-heredity:vocabulary}{等位基因}的标准而言，这是极强的选择。
\textbf{6.} 设 $w_{MM} = w_{Mm} = 1 - s$，$w_{mm} = 1$，则 $\bar w = 1
- s(1 - q^{2})$，$p' = p(1 - s)/\bar w$，因此 $\Delta p = p[(1 - s) -
\bar w]/\bar w = -spq^{2}/\bar w$。淡色型仅以
\omterm{def:b2:meiosis-heredity:vocabulary}{纯合}子 $mm$ 的形式受到青睐，开始时其比例为 $q^{2} = 0.1$，因此
回归起初很缓慢。当 $q \to 1$ 时，$\Delta p \approx -sp$，$M$
呈指数下降，每代乘以因子 $1 - s$：\omterm{def:b2:meiosis-heredity:vocabulary}{显性}\omterm{def:b2:meiosis-heredity:vocabulary}{等位基因}
无处可藏。上升过程恰为其镜像：
$\Delta p \approx +sp$ 当 $M$ 稀少时（指数式起步），随后
$\Delta q \approx -sq^{2}p \to$ 进展缓慢，淡色\omterm{def:b2:meiosis-heredity:vocabulary}{等位基因}隐藏于
\omterm{def:b2:meiosis-heredity:vocabulary}{杂合}子中。递推式给出回归所需的 27 代，
即从 $p = 0.68$ 降至 $0.001$，其中 16 代用于使其降至 \qty{5}{\%}。
\textbf{7.} $24$ 个基因拷贝：$0.95^{24} = 0.29$。
\textbf{8.} $H_1 = 0.4\,(1 - 1/24) = 0.383$。
\textbf{9.} $H_{100} = 0.383\,(1 - 1/80)^{100} = 0.383\times 0.284 =
0.109$：为大陆 $0.4$ 的 \qty{27}{\%}。
\textbf{10.} $P(\text{固定}) = 1/(2N) = 1/80 = 0.0125$；期望
固定所需时间为 $4N = 160$ 代。
\textbf{11.} 无\omterm{prop:b2:population-genetics:inbreeding}{近交}时 $q^{2} = 0.0025$；当 $F = 0.3$ 时，
$q^{2} + Fpq = 0.0025 + 0.3\times 0.95\times 0.05 = 0.017$，
患病出生者为原来的七倍。
\textbf{12.} 每代一个迁入个体（$m = 1/40$）即足以
防止岛屿在中性位点上发生\omterm{def:b2:cell-differentiation:differentiation}{分化}，并
补充漂变所丢失的\omterm{def:b2:meiosis-heredity:vocabulary}{等位基因}；它还输入了
大陆的有害\omterm{def:b2:meiosis-heredity:vocabulary}{等位基因}，其频率与大陆相同，这些\omterm{def:b2:meiosis-heredity:vocabulary}{等位基因}在
岛上\omterm{prop:b2:population-genetics:inbreeding}{近交}的条件下以\omterm{def:b2:meiosis-heredity:vocabulary}{纯合}子形式暴露出来——
其\omterm{thm:b2:population-genetics:balance}{遗传负荷}由迁入与暴露之间的平衡决定。
\textbf{13.} $\hat q = \sqrt{10^{-5}} = 0.0032$；患病出生者 $\hat
q^{2} = 10^{-5}$。
\textbf{14.} $\hat q = \sqrt{10^{-5}/0.1} = 0.01$；患病出生者
$10^{-4}$：选择强度减弱为十分之一，
使该\omterm{def:b2:population-genetics:frequencies}{等位基因频率}变为原来的约三倍，疾病则为原来的十倍。
\textbf{15.} $\hat p = \mu/(hs) = 10^{-5}/0.05 = 2\times 10^{-4}$；
携带者 $2\hat p = 4\times 10^{-4}$。
\textbf{16.} 从 $\hat q = 0.0032$ 升至 $0.01$：该\omterm{def:b2:meiosis-heredity:vocabulary}{等位基因}以由
\omterm{def:b2:genome-diversification:mutation}{突变}决定的速率上升，即 $\mu = 10^{-5}$（每代），因此
接近该值需约 $\Delta q/\mu \approx 700$ 代——
约 \num{20000} 年：如今作出的任何决定，都将影响一个
与我们今天截然不同的种群。
\textbf{17.} 每个位点贡献 $2\hat q = 0.0063$ 个致死拷贝
每人；在 \num{20000} 个位点上，约为 $130$ 个——但这一假设
是每个基因都能\omterm{def:b2:genome-diversification:mutation}{突变}为\omterm{def:b2:meiosis-heredity:vocabulary}{隐性}致死；若采用常用估计
的必需基因约有数千个，一个人在\omterm{def:b2:meiosis-heredity:vocabulary}{杂合}状态下携带
数个至数十个致死当量。
\textbf{18.} 每个致死拷贝位于不同位点，且均罕见，因此
一对随机配偶同时携带同一个的概率很小；
表兄妹通过共同祖先共享八分之一的基因，因此对于
表兄妹各自携带的若干致死拷贝中的每一个，子女为\omterm{def:b2:meiosis-heredity:vocabulary}{纯合}子的风险为
每个共享致死拷贝 $1/16$——表兄妹婚配子女的超额死亡率
经测量约为百分之几。
\textbf{19.} $S = 1.16\times \qty{800}{L} = \qty{930}{L}$；$R = 0.4\times
930 = \qty{370}{L}$。
\textbf{20.} 五代后约为 \qty{1860}{L}；实际上会更小，
因加性方差已耗尽，且 $h^{2}$
是在原始畜群中估计的，也因为
使顶尖母牛表现优异的环境（饲料、圈舍）不一定能传递下去。
\textbf{21.} 平均选择差为 $(1.16 + 2.4)/2 = 1.78$
个标准差，$S = \qty{1424}{L}$，$R = 0.4\times 1424 = \qty{570}{L}$。
\textbf{22.} \omterm{def:b2:asexual-reproduction:clone}{克隆}个体没有遗传方差，因此 $h^{2} = 0$，且 $R =
0$，无论 $S$ 为何值：\omterm{thm:b2:population-genetics:heritability}{遗传力}是种群的属性，即其
方差中属于遗传的比例，而不是性状本身的属性。
\textbf{23.} $R = 0.7\times 0.6 = \qty{0.42}{mm}$，即更深。
\textbf{24.} 获得同样的选择响应 $\qty{370}{L}$ 需要 $S = 370/0.1 =
\qty{3700}{L}$，即 $4.6$ 个标准差——只留下
不到万分之一的母牛，这并不可行，因此
低\omterm{thm:b2:population-genetics:heritability}{遗传力}性状通常改良缓慢，或借助其他手段
（如后代测定、扩大家系）。
\textbf{25.} 蛾：$s \approx 0.2$ 至 $0.3$，在 $0.3$ 时为 28 代，
在 $0.2$ 时为 45 代，与观测的 50 代相比；岛屿：$H_{100} \approx 0.11$，
为大陆的 \qty{27}{\%}；畜群：仅由母牛每代提供 \qty{370}{L}，
若选用优良公牛则为 \qty{570}{L}。
\end{solution}
